\section{Regional normalization and filtered vectors}

The original nearest-neighbour chain admits the scalar normalization used in
\cite[Section~6.1]{Arad2013AreaLaw}. Separately, supplied approximate projectors
improve normalized low-rank seeds with explicit phase and error estimates.
These results do not assert existence of a favorable approximate projector.

\begin{lemma}[Disjoint regional ground-energy shifts]
    \label{lem:aklv_disjoint_ground_shifts}
    \lean{TNLean.MPS.AreaLaw.AKLV.exists_common_ground_shifts_of_disjoint_support,
        TNLean.MPS.AreaLaw.AKLV.exists_three_disjoint_ground_shifts}
    \leanok
    Two Hermitian operators on disjoint site regions admit scalar shifts making
    both positive, with a common normalized vector attaining both shifted zero
    energies. Three pairwise disjoint regions admit positive shifts and a
    normalized zero-energy test for their sum. Each shift is nonnegative if
    its original operator is positive.
\end{lemma}
\begin{proof}\leanok
    In cut coordinates the two operators are $M\otimes I$ and $I\otimes N$.
    Tensor minimum-energy eigenvectors and transport back to the original
    sites. For three regions, first normalize two and then their sum against
    the third. Attained lower energies add. Positivity of an original operator
    makes its attained minimum nonnegative.
\end{proof}

\begin{lemma}[Normalization of the original chain]
    \label{lem:aklv_original_chain_normalization}
    \lean{TNLean.MPS.AreaLaw.AKLV.exists_chain_normalization,
        TNLean.MPS.AreaLaw.AKLV.middleCutBlock_le_window}
    \leanok
    Let $0\le h_i\le I$ be nearest-neighbour interactions with a unit ground
    vector and nonnegative projector-gap lower bound. Removing two distinct
    bonds gives three regions. There are nonnegative scalar shifts making
    their three blocks positive, preserving the original ground vector and
    gap after the total scalar shift, and giving ground energy $0\le e\le2$.
    The shifted middle block is bounded above by the distance between the
    two removed bonds times the identity.
\end{lemma}
\begin{proof}\leanok
    \uses{lem:aklv_disjoint_ground_shifts}
    The three bond sums have disjoint supports, and their shifted sum has a
    normalized zero-energy test. Restoring two positive contractions adds at
    most two to its energy. The variational principle gives the stated ground
    energy bound. A nonnegative middle shift cannot increase its original sum,
    which is bounded by its number of unit-bounded bonds.
\end{proof}

\begin{lemma}[Normalized orthogonal decomposition]
    \label{lem:aklv_phase_normalization}
    \lean{TNLean.MPS.AreaLaw.AKLV.inner_sub_inner_smul_unit_eq_zero,
        TNLean.MPS.AreaLaw.AKLV.exists_phase_normalized_of_decomposition,
        TNLean.MPS.AreaLaw.AKLV.schmidtRank_phase_normalized_eq}
    \leanok
    If $\chi=a\phi+v$, $\|\phi\|=1$, $a\ne0$ and $v\perp\phi$, then $\chi\ne0$
    and some unit phase makes its normalization have squared distance at most
    $2\|v\|^2/|a|^2$ from $\phi$. The normalization and phase preserve Schmidt rank.
    In particular, subtracting $\langle\phi,\psi\rangle\phi$ from a vector
    leaves a vector orthogonal to the unit vector $\phi$.
\end{lemma}
\begin{proof}\leanok
    Pythagoras gives $\|\chi\|^2=|a|^2+\|v\|^2$ and a nonzero denominator.
    The normalized orthogonal mass is at most $\|v\|^2/|a|^2$. Minimizing
    over phase bounds squared distance by twice this mass. Both scalar factors
    are nonzero, so they preserve the coefficient matrix rank.
\end{proof}

\begin{lemma}[Normalized filtering near a fixed target]
    \label{lem:aklv_nearby_filtered_vector}
    \lean{TNLean.MPS.AreaLaw.AKLV.exists_normalized_approximation_step}
    \leanok
    Let $K$ fix a unit vector $\phi$, preserve its orthogonal complement, and
    have squared shrinking $\delta\ge0$. If a unit vector $\psi$ satisfies
    $\|\phi-\psi\|^2\le r\le1/2$, then $K\psi\ne0$ and a phase-corrected
    normalization has squared distance at most $4\delta r$ from $\phi$.
\end{lemma}
\begin{proof}\leanok
    \uses{lem:aklv_phase_normalization}
    The squared target component is at least $1/2$ and the input orthogonal
    mass is at most $r$. Filtering reduces that mass by $\delta$; normalize.
\end{proof}

\begin{lemma}[Filtering a seed with nonzero overlap]
    \label{lem:aklv_filtered_seed}
    \lean{TNLean.MPS.AreaLaw.AKLV.exists_filtered_approximation_of_overlap,
        TNLean.MPS.AreaLaw.AKLV.exists_filtered_approximation_of_seven_eighths}
    \leanok
    Suppose $K$ fixes a unit target, preserves its orthogonal complement with
    squared shrinking $\delta\ge0$, and has operator Schmidt rank at most $D$.
    If a unit seed has Schmidt rank at most $R$ and overlap at least $\mu>0$,
    its filtered output is nonzero and admits a phase-corrected normalization
    with squared error at most $2\delta/\mu^2$ and Schmidt rank at most $DR$.
    For $\mu=7/8$ and $\delta\le(49/128)e^{-2x}$ its norm error is at most $e^{-x}$.
\end{lemma}
\begin{proof}\leanok
    \uses{lem:aklv_phase_normalization}
    The seed's orthogonal squared norm is at most one. Filtering reduces it by
    $\delta$ while preserving squared target component at least $\mu^2$.
    Normalize and apply the operator-to-vector Schmidt-rank bound.
\end{proof}
